% Formulario de consentimiento del conductor y del cobrador, para imprimir
% y firmar. Dice LO MISMO que /consentimiento (server/consentimiento.html) en
% su versión 2026-11-11: si una cambia, cambia la otra, y se sube
% CONSENTIMIENTO_VERSION en server/index.js.
\documentclass[10.5pt,a4paper]{article}
% Preámbulo común de los documentos legales (legal/*.tex).
% Se compila cada uno con:   cd legal && latexmk -pdf <archivo>.tex
% La paleta es la del producto y la del manual: si cambia una, cambian las dos.
\usepackage[T1]{fontenc}
\usepackage[utf8]{inputenc}
\usepackage[spanish,es-noshorthands,es-tabla]{babel}
\usepackage[scaled=0.92]{helvet}
\renewcommand{\familydefault}{\sfdefault}
\usepackage{amssymb}
% Los que usan los textos y pdflatex no trae
\DeclareUnicodeCharacter{2713}{\checkmark}
\DeclareUnicodeCharacter{22EF}{\ensuremath{\cdots}}
\DeclareUnicodeCharacter{2192}{\ensuremath{\rightarrow}}
\usepackage[a4paper,top=2cm,bottom=2.2cm,left=2.2cm,right=2.2cm,headheight=14pt]{geometry}
\usepackage[table]{xcolor}
\usepackage{microtype}
\usepackage{enumitem}
\usepackage{booktabs}
\usepackage{tabularx}
\usepackage{array}
\usepackage{longtable}
\usepackage{fancyhdr}
\usepackage{titlesec}
\usepackage{lastpage}
\usepackage[most]{tcolorbox}
\usepackage[hidelinks]{hyperref}

\definecolor{tinta}{HTML}{0B2239}
\definecolor{tinta2}{HTML}{2A4864}
\definecolor{tenue}{HTML}{496A86}
\definecolor{marca}{HTML}{1668C4}
\definecolor{ambar}{HTML}{A67300}
\definecolor{rojo}{HTML}{C2001D}
\definecolor{fondo}{HTML}{F4F7FB}
\definecolor{linea}{HTML}{C3D4E4}
\definecolor{ambarfondo}{HTML}{FDF6E4}
\definecolor{ambarborde}{HTML}{E4C88A}
\definecolor{azulfondo}{HTML}{E7F0FB}
\definecolor{azulborde}{HTML}{C9DDF4}

\color{tinta}
\setlength{\parindent}{0pt}
\setlength{\parskip}{5pt plus 2pt}
\linespread{1.08}
\setlist{itemsep=2pt,topsep=3pt}

\titleformat{\section}{\color{tinta}\large\bfseries}{}{0em}{}
\titleformat{\subsection}{\color{tinta2}\normalsize\bfseries}{}{0em}{}
\titlespacing*{\section}{0pt}{12pt}{4pt}
\titlespacing*{\subsection}{0pt}{8pt}{2pt}

\renewcommand{\arraystretch}{1.25}
\arrayrulecolor{linea}

% Lo que hay que completar a mano: una línea para escribir
\newcommand{\llenar}[1][6cm]{\rule[-2pt]{#1}{0.5pt}}
% Lo que tiene que decidir o completar el abogado o las partes
\newcommand{\completar}[1]{{\color{ambar}[#1]}}
% Una casilla para marcar
\newcommand{\casilla}{\fbox{\rule{0pt}{7pt}\rule{7pt}{0pt}}}

\newtcolorbox{nota}[1][]{colback=azulfondo,colframe=azulborde,boxrule=0.6pt,arc=3pt,
  left=8pt,right=8pt,top=5pt,bottom=5pt,fontupper=\small\color{tinta2},#1}
\newtcolorbox{borrador}{colback=ambarfondo,colframe=ambarborde,boxrule=0.6pt,arc=3pt,
  left=8pt,right=8pt,top=5pt,bottom=5pt,fontupper=\small}

% Pie: qué documento es, su versión y la página
\newcommand{\documento}[2]{%
  \pagestyle{fancy}\fancyhf{}%
  \renewcommand{\headrulewidth}{0pt}%
  \fancyfoot[L]{\scriptsize\color{tenue}#1 · #2}%
  \fancyfoot[R]{\scriptsize\color{tenue}Página \thepage\ de \pageref*{LastPage}}%
}

% Firma: una línea y lo que va debajo
\newcommand{\firma}[1]{\begin{minipage}[t]{0.44\textwidth}\vspace{1.6cm}\rule{\linewidth}{0.5pt}\\{\small\color{tinta2}#1}\end{minipage}}

\begin{document}
\documento{Consentimiento para el tratamiento de datos personales}{versión 2026-11-11}

{\LARGE\bfseries Consentimiento para el tratamiento de datos personales}\par
{\small\color{tenue}Aplicación «Control de ruta» · versión 2026-11-11 · Ley N.º 29733 y su Reglamento}

\begin{borrador}
\textbf{Borrador.} Lo tiene que revisar el abogado antes de usarlo
(ver \texttt{legal/LEEME.md}). Se firma uno por persona: la cooperativa guarda
el original y le entrega una copia. Después, en Despacho → Gestión →
Personas → \emph{Consentimiento}, se anota la fecha de la firma.
\end{borrador}

\section{Quién trata sus datos}
\textbf{Responsable:} la cooperativa \llenar[7.5cm], RUC \llenar[3.2cm],
con domicilio en \llenar[9.5cm]. Ella decide para qué se usan sus datos y
es a quien usted les reclama.

\textbf{Encargado:} el operador del sistema de control de flota que la
cooperativa contrató, que guarda y procesa los datos solo por cuenta de ella
y según sus instrucciones. Contacto: \llenar[8cm].

\section{Qué datos y para qué}
\begin{itemize}[leftmargin=14pt]
  \item \textbf{Su identidad en el sistema:} nombre, alias, la combi que maneja y su rol.
  \item \textbf{Su ubicación durante el turno}, de forma continua,
    \textbf{también con la pantalla apagada o la aplicación en segundo plano}.
    Sirve para mostrar la combi en el mapa de la cooperativa, medir la
    separación con las demás, registrar las vueltas, los desvíos y las horas
    en ruta con las que la cooperativa liquida el trabajo, y ubicarlo si pide
    ayuda (SOS). Fuera del turno no se envía.
  \item \textbf{En la pantalla de inicio de turno}, con la aplicación abierta,
    el teléfono mira la ubicación \textbf{sin enviarla} para darse cuenta de
    si usted ya está manejando sin haber salido a ruta: le avisa y, si no
    contesta en un minuto, sale a ruta por usted.
  \item \textbf{Lugares puntuales del turno}, con su hora: dónde se declaró
    en ruta, ausente o fuera ---incluido \textbf{dónde terminó el turno}---,
    dónde se cortó la señal, dónde estuvo parado y dónde salió del trazado.
    No es el recorrido completo: son esos puntos.
  \item \textbf{Marcas automáticas:} el sistema anota por su cuenta señales
    como «GPS sospechoso», «entrada tardía», «salió solo» o «fuera de ruta».
    Son avisos para que una persona de la cooperativa los mire:
    \textbf{el sistema no sanciona ni decide nada por sí solo}, y usted puede
    pedir que se revisen y que se le explique cualquiera.
  \item \textbf{El chat de la ruta:} los textos, notas de voz y fotos que
    usted decide mandar, y las emergencias (SOS).
  \item \textbf{El estado del teléfono} mientras está en ruta (batería,
    ahorro de batería, versión de la app): solo el último valor, no se guarda.
\end{itemize}
No se usan para publicidad, no se venden y no se entregan a terceros para fines propios.

\section{Quién los ve y dónde se guardan}
El personal de Despacho y la gerencia de \textbf{su} cooperativa, y el
operador del sistema para darle soporte. Para funcionar, el sistema usa
proveedores de alojamiento y de mapas cuyos \textbf{servidores están fuera
del Perú}; las copias de respaldo se guardan cifradas. La lista completa
está en la política de privacidad del sistema (\texttt{/privacidad}).

\section{Cuánto tiempo}
\begin{tabularx}{\linewidth}{>{\raggedright}p{6.4cm}X}
\toprule
Chat (texto, fotos y notas de voz) & 30 días; lo que Despacho guarda como evidencia de un hecho, hasta 365 días\\
Vueltas y desvíos & 120 días\\
Lugares puntuales del turno y marcas automáticas & 120 días\\
Turnos (horas trabajadas) y emergencias (SOS) & 365 días\\
Registro de actividad & inicios de sesión 90 días; lo del recorrido 365 días; altas, bajas y cambios de clave 2 años\\
Posición en vivo & solo la última conocida (de lo demás, sólo los lugares puntuales de arriba)\\
\bottomrule
\end{tabularx}

\section{Sus derechos}
Puede pedir \textbf{acceso} a sus datos, su \textbf{rectificación}, su
\textbf{cancelación}, \textbf{oponerse} a su uso y pedir su
\textbf{portabilidad}, y puede \textbf{revocar este consentimiento} cuando
quiera, sin efecto hacia atrás. Se pide a la cooperativa (a su central) o al
contacto del encargado. Si no le responden, puede reclamar ante la Autoridad
Nacional de Protección de Datos Personales.

La ubicación durante el turno es la función central de la aplicación: sin
ella no puede usarse para trabajar. La cooperativa le dirá cómo se registra
su trabajo si no da o revoca este consentimiento.

\begin{samepage}
\section{Lo que usted autoriza}
Marque cada casilla que acepta. Cada una es independiente.

\begin{tabularx}{\linewidth}{@{}p{14pt}X@{}}
\casilla & Acepto que se use mi \textbf{ubicación durante el turno, también con la pantalla apagada o la aplicación en segundo plano}, para los fines de arriba.\\
\casilla & Acepto que, en la pantalla de inicio de turno, la aplicación mire mi ubicación sin enviarla y \textbf{salga a ruta por mí} si ve que estoy manejando y no contesto.\\
\casilla & Acepto que mis datos se guarden en \textbf{servidores fuera del Perú}, con las protecciones descritas.\\
\casilla & Declaro que recibí una copia de este documento y que puedo leer la política de privacidad completa en \texttt{/privacidad}.\\
\end{tabularx}

\vspace{10pt}
Nombre completo: \llenar[11cm]\par\vspace{6pt}
DNI: \llenar[4cm]\quad Usuario en la app: \llenar[5cm]\par\vspace{6pt}
Lugar y fecha: \llenar[10cm]

\noindent\firma{Firma del conductor o cobrador}\hfill\firma{Por la cooperativa (nombre, cargo y firma)}
\end{samepage}
\end{document}
